%=====================================================================
%  Aerotrack: Sistema De Estimacion De Aglomeraciones y Tracking De
%  Personas Mediante Computer Vision
%  Universidad ESAN -- Facultad de Ingenieria
%  Carrera de Ingenieria de Tecnologias de Informacion y Sistemas
%
%  Compilacion en Overleaf: pdfLaTeX (motor por defecto), dos pasadas.
%  No requiere biber/biblatex: las referencias se listan como texto
%  con formato APA en secciones/referencias.tex.
%
%  Archivos .tex codificados en UTF-8.
%=====================================================================
\documentclass[12pt,a4paper]{article}

%--------------------- Codificacion, fuente e idioma ------------------
\usepackage[T1]{fontenc}
\usepackage[utf8]{inputenc}
% Times New Roman 12 en texto y formulas.
\usepackage{newtxtext}
\usepackage{newtxmath}
\usepackage[spanish,es-tabla,es-noquoting,es-lcroman]{babel}
\usepackage{csquotes}

%--------------------- Geometria y estilo ----------------------------
\usepackage[a4paper,left=3cm,right=2.5cm,top=3cm,bottom=2.5cm,headheight=15pt]{geometry}
\usepackage{fancyhdr}
\usepackage{setspace}
\usepackage{indentfirst}
\setlength{\parindent}{1.25cm}
\setlength{\parskip}{0pt}

%--------------------- Tablas, listas, graficos ----------------------
\usepackage{graphicx}
\usepackage{float}
\usepackage{xcolor}
\usepackage{booktabs}
\usepackage{array}
\usepackage{tabularx}
\usepackage{enumitem}
\usepackage{ragged2e}
\usepackage{longtable}
\usepackage{tikz}
\usetikzlibrary{positioning,arrows.meta,shapes.geometric,calc}
% Permite cortar las direcciones web largas de las referencias en los guiones.
\PassOptionsToPackage{hyphens}{url}
\usepackage{hyperref}
\Urlmuskip=0mu plus 1mu

% Todo el documento va en negro: texto, tablas, diagramas y enlaces.
\color{black}

% Tipos de columna para tablas anchas (casos de uso, requisitos).
\newcolumntype{L}[1]{>{\RaggedRight\arraybackslash}p{#1}}
\newcolumntype{C}[1]{>{\centering\arraybackslash}p{#1}}
\newcolumntype{Y}{>{\RaggedRight\arraybackslash}X}

% Tablas: booktabs sin rayas verticales ni sombreado, con un poco de aire.
\renewcommand{\arraystretch}{1.15}

%--------------------- Figuras ----------------------------------------
% Todas las capturas y figuras se centran al mismo ancho.
\newlength{\anchofigura}
\setlength{\anchofigura}{0.85\textwidth}
%   \figura{archivo}{pie}{etiqueta}
\newcommand{\figura}[3]{%
  \begin{figure}[H]
    \centering
    \includegraphics[width=\anchofigura]{imagenes/#1}
    \caption{#2}
    \label{#3}
  \end{figure}}

%--------------------- Marcadores [VERIFICAR] -------------------------
% \verificar{texto} marca una afirmacion que sigue sin confirmar contra
% su fuente. NO se imprime en el PDF: el texto queda en el fuente y cada
% marcador deja rastro en el registro de compilacion, con un recuento
% final ("MARCADORES [VERIFICAR] EN EL CUERPO: n").
% La macro abre con \unskip para no dejar doble espacio cuando el
% marcador va en una linea propia, y no termina en \ignorespaces para no
% comerse el espacio siguiente.
\newcounter{verificarcount}
\newcommand{\verificar}[1]{%
  \ifhmode\unskip\fi
  \stepcounter{verificarcount}%
  \typeout{[VERIFICAR \theverificarcount] marcador oculto en pagina \thepage.}}
\AtEndDocument{%
  \typeout{^^J=====================================================^^J%
    MARCADORES [VERIFICAR] EN EL CUERPO: \theverificarcount^^J%
    =====================================================^^J}}

%--------------------- Hiperenlaces -----------------------------------
\hypersetup{
  unicode=true,
  hidelinks,
  pdftitle={Aerotrack: Sistema de Estimacion de Aglomeraciones y Tracking
            de Personas Mediante Computer Vision},
  pdfauthor={Aguilar Contreras, Angel Jesus; Rivadeneyra Huaman, Stephano;
             Ortega, Jose; Moreno Ugarte, Fabian},
  pdfsubject={Universidad ESAN, Facultad de Ingenieria, Carrera de
              Ingenieria de Tecnologias de Informacion y Sistemas}
}

%--------------------- Encabezado y pie --------------------------------
% Sin filetes: ni bajo el encabezado ni sobre el pie.
\pagestyle{fancy}
\fancyhf{}
\fancyhead[L]{\small Aerotrack}
\fancyhead[R]{\small Universidad ESAN}
\fancyfoot[C]{\small\thepage}
\renewcommand{\headrulewidth}{0pt}
\renewcommand{\footrulewidth}{0pt}
\fancypagestyle{plain}{%
  \fancyhf{}%
  \fancyhead[L]{\small Aerotrack}%
  \fancyhead[R]{\small Universidad ESAN}%
  \fancyfoot[C]{\small\thepage}%
  \renewcommand{\headrulewidth}{0pt}%
  \renewcommand{\footrulewidth}{0pt}}

\onehalfspacing

%=====================================================================
\begin{document}
%=====================================================================

%--------------------------- PORTADA ---------------------------------
\begin{titlepage}
\centering
\singlespacing
\vspace*{0.2cm}

\IfFileExists{imagenes/logo-esan.png}{%
  \includegraphics[width=5cm]{imagenes/logo-esan.png}%
}{%
  \framebox[5cm][c]{\rule{0pt}{1.6cm}\small\itshape Logo ESAN}%
}

\vspace{1.2cm}

{\LARGE\bfseries Aerotrack: Sistema De Estimación De Aglomeraciones y\par}
\vspace{0.2cm}
{\LARGE\bfseries Tracking De Personas Mediante Computer Vision\par}

\vspace{1.4cm}

{\large Universidad ESAN.\par}
\vspace{0.2cm}
{\large Facultad de Ingeniería.\par}
\vspace{0.2cm}
{\large Carrera de Ingeniería de Tecnologías de Información y Sistemas.\par}

\vspace{0.6cm}
{\large Docente: Fabian Arteaga, Junior.\par}

\vspace{0.9cm}

\begin{tabular}{@{}L{5.5cm}L{2.3cm}L{6cm}@{}}
\toprule
\textbf{Integrante} & \textbf{Código} & \textbf{Rol} \\
\midrule
Aguilar Contreras, Angel Jesus & 22200133 & Liderazgo e integración \\[0.15cm]
Rivadeneyra Huaman, Stephano & 22101822 & Ciencia de datos \\[0.15cm]
Ortega, José & 25200719 & Ingeniería de datos \\[0.15cm]
Moreno Ugarte, Fabián & 25200717 & Validación ética, legal y documentación \\
\bottomrule
\end{tabular}

\vfill

{\large 6 de octubre de 2026.\par}
\vspace{0.3cm}
{\large Lima, Perú.\par}


\end{titlepage}

%--------------------------- INDICE ----------------------------------
\pagenumbering{roman}
\setcounter{tocdepth}{2}
{\singlespacing\tableofcontents}
\clearpage
{\singlespacing\listoffigures}
\clearpage
{\singlespacing\listoftables}
\clearpage
\pagenumbering{arabic}

%------------------------- SECCIONES ---------------------------------
\section{Resumen Ejecutivo}
\label{sec:resumen-ejecutivo}

Este informe documenta AeroTrack, un sistema de visión computacional que estima la
concentración de personas por zona y sigue su recorrido entre cámaras sobre el plano de un
espacio, pensado para el Aeropuerto Internacional Jorge Chávez a partir de sus cámaras fijas.
El sistema no recurre a reconocimiento facial, pero sí individualiza a las personas: les asigna un
identificador anónimo, las sigue dentro de cada cámara y las vuelve a reconocer entre cámaras
por la apariencia de su cuerpo. Esa capacidad sitúa al proyecto de lleno bajo la Ley 29733 de
Protección de Datos Personales y su reglamento vigente (Decreto Supremo 016-2024-JUS), y el
informe la documenta tal como está implementada, sin presentarla como un tratamiento
agregado que no es.

El marco teórico revisa los trabajos que sostienen el prototipo. CSRNet, P2PNet y P2R se
estudian como antecedentes del conteo de multitudes, que orientaron la discusión sobre la
granularidad de la salida, aunque el prototipo final no usa ninguno de los tres. El sistema
construido descansa en tres piezas: un detector de personas YOLO de la familia YOLO11
(Jocher y Qiu, 2024), cuya variante se elige según el hardware disponible; un seguimiento por
cámara que aplica la asociación en dos etapas de ByteTrack (Zhang et al., 2022) sobre el punto
de los pies, reforzada con señales de BoT-SORT (Aharon et al., 2022); y un descriptor de
apariencia OSNet (Zhou et al., 2019) que sostiene la reidentificación entre cámaras.

El prototipo funciona de punta a punta en un equipo sin GPU dedicada: configura el plano y
las cámaras, calibra cada cámara con una homografía, detecta y sigue a las personas, las asocia
entre cámaras, emite alertas de aglomeración con aviso por correo y produce reportes de
ocupación, flujo y oportunidad comercial. Su evaluación disponible es limitada y así se declara:
la identidad entre cámaras se midió sobre dos cámaras de demostración con cuatro personas y
veinticuatro segundos de video, con umbrales elegidos sobre esos mismos datos, de modo que
sus cifras son una demostración de funcionamiento y no una medida de precisión general.

La especificación de requerimientos adopta la numeración maestra del equipo: dieciocho casos
de uso, dieciocho requisitos funcionales y ocho no funcionales, con su matriz de trazabilidad.
El requisito no funcional rector, RNF-01, prohíbe derivar del video rasgos que permitan
identificación biométrica o reidentificación individual persistente. Con la configuración por
defecto del prototipo \textbf{ese requisito no se cumple}: la memoria de apariencia, activa salvo
que se desactive, guarda en el disco del equipo el vector OSNet de cuerpo entero de cada
persona durante veinticuatro horas (configurable hasta siete días) y la reconoce entre sesiones y
entre días. La corrección está identificada y depende solo del equipo: desactivar esa memoria
por defecto, de modo que el vector viva únicamente en memoria durante la sesión. A ese
hallazgo se suma el incumplimiento de la medida de salida agregada (M-03), porque las
grabaciones de cada sesión conservan posiciones e identificadores por persona sin un plazo de
retención automático.

Las recomendaciones plantean cerrar primero esos dos hallazgos, formalizar la Evaluación de
Impacto en Protección de Datos cuyo contenido sustantivo recoge este informe, validar la
exactitud del sistema con material del propio terminal y, solo después, considerar el despliegue
en infraestructura de nube para escalar a más cámaras simultáneas y la integración con fuentes
operativas de los locales del terminal. Ambas extensiones dependen de la autorización de LAP,
y la primera exige analizar antes el flujo transfronterizo de datos que supondría procesar fuera
del Perú (Anexo~\ref{anx:marco}).

\clearpage
\section{Introducción}
\label{sec:introduccion}

\subsection{Contextualización del Problema}
\label{subsec:contextualizacion}

El Aeropuerto Internacional Jorge Chávez, operado por Lima Airport Partners (LAP), no
distribuye a sus pasajeros de manera uniforme a lo largo del día. En determinadas franjas
horarias y en puntos concretos del terminal (filtros de seguridad, control migratorio,
mostradores de \emph{check-in}, salas de embarque) se forman concentraciones de personas
notablemente superiores a las del resto del recinto en ese mismo momento. La aglomeración,
por sí sola, no es el problema: lo es cuando la cantidad de personas reunidas en un punto
excede, de forma sostenida, la capacidad de atención dispuesta para ese punto, y aparecen
colas que se alargan, tiempos de espera que comprometen la conexión de un vuelo, o
condiciones de incomodidad en un espacio cerrado. LAP cuenta ya con infraestructura de
videovigilancia de circuito cerrado en el terminal; este proyecto explora si esa misma
infraestructura puede aprovecharse, mediante visión computacional, para producir una
estimación automatizada de esa concentración y de los recorridos entre cámaras, sin
reconocimiento facial. Conviene adelantar desde aquí lo que el resto del informe desarrolla: para
seguir a una persona de una cámara a otra el sistema la individualiza con un identificador
anónimo y la reconoce por la apariencia de su cuerpo, y esa decisión de diseño es la que fija su
régimen jurídico (Secciones~\ref{subsec:restricciones} y~\ref{subsec:analisis-teas}).

El equipo no cuenta, al momento de escribir este informe, con evidencia documentada por LAP (incidentes de
aglomeración registrados, tiempos de espera medidos, quejas de pasajeros contabilizadas) que
cuantifique el problema en el propio terminal.\verificar{incidentes por aglomeración
documentados, tiempos de cola observados y quejas recibidas. Solicitado a LAP como punto C-07.} El proyecto parte, por tanto, de una necesidad
operativa plausible y compartida por la literatura de conteo de multitudes en espacios de alto
tránsito, y no de una cifra propia de incidentes del Aeropuerto Jorge Chávez.

\subsection{Importancia e Impacto}
\label{subsec:importancia-impacto}

Conocer cuándo y dónde se concentran las personas es información operativa de valor
directo para un operador aeroportuario: permite anticipar la apertura de mostradores
adicionales, redirigir el flujo de pasajeros hacia un filtro descargado, o dimensionar al personal
de atención con mayor precisión que una revisión manual y periódica, que cubre un punto a la
vez y no deja una serie homogénea sobre la que puedan compararse dos días, dos zonas o dos
franjas horarias. El impacto del proyecto tiene, además, una dimensión que excede lo
operativo: un sistema que observa continuamente a personas en un espacio de uso público
concesionado activa, por diseño, un examen de proporcionalidad bajo la Ley 29733, y el
análisis de impacto de este informe documenta un riesgo distributivo real: si el sistema
subcuenta sistemáticamente a un grupo de personas, las zonas donde ese grupo se concentra
reciben menos recursos de forma sistemática, un daño indirecto pero real que no depende de
identificar a ninguna persona en particular.

\subsection{Motivación desde la Ingeniería}
\label{subsec:motivacion-ingenieria}

Desde la perspectiva de Ingeniería de Sistemas, el proyecto es un caso de estudio útil por
una razón poco común: la restricción dominante no es de cómputo ni de datos, sino normativa.
La mayoría de los trabajos de conteo de multitudes revisados en el marco teórico optimizan
exactitud bajo el supuesto implícito de que capturar más información por persona siempre
mejora el sistema; este proyecto invierte esa prioridad, porque el principio de proporcionalidad
de la Ley 29733 (art. 7) exige que el tratamiento sea adecuado, relevante y no excesivo respecto
de su finalidad. Diseñar bajo esa restricción obliga a hacer explícitas decisiones que en otros
contextos quedarían implícitas: qué se persiste, por cuánto tiempo, y qué mecanismo de
desempate se usa cuando dos lecturas del sistema son ambiguas. La segunda motivación es de
integración bajo restricciones reales de hardware: el prototipo se ejecuta en CPU, sin GPU
dedicada, en computadoras portátiles universitarias, combinando piezas de procedencia distinta
(un detector de personas preentrenado, un seguimiento adaptado, un descriptor de apariencia
preentrenado, una homografía calibrada a mano, una base de datos local y una consola web) en
un único flujo coherente. El costo de esa restricción es concreto: en un procesador de gama
media sin GPU, analizar varias cámaras con una muestra cada 0,2 segundos va entre cinco y
veinticinco veces más lento que el video, por lo que el sistema ofrece modos de rendimiento que
espacian las muestras según el número de cámaras.\verificar{rango de 5 a 25 veces más lento que
el video: estimación del equipo sobre un procesador i5 sin GPU, no una medición formal con
protocolo documentado.}

\subsection{Objetivo General}
\label{subsec:objetivo-general}

Desarrollar y evaluar un sistema de visión computacional capaz de estimar la
concentración y tracking de personas en distintas zonas del Aeropuerto Internacional Jorge
Chávez a partir de imágenes de cámaras fijas, sin recurrir a reconocimiento facial ni a la
extracción de rasgos biométricos identificables.

\subsection{Objetivos Específicos}
\label{subsec:objetivos-especificos}

\begin{enumerate}
  \item Revisar el estado del arte en detección, localización y conteo de personas en escenas
  congestionadas, e identificar qué enfoque técnico resulta más adecuado para un entorno
  aeroportuario y para las restricciones de privacidad del proyecto.
  \item Determinar el marco de cumplimiento normativo aplicable al tratamiento de imágenes
  de un espacio de uso público bajo gestión concesionada, incluyendo el análisis de
  impacto y los requisitos de privacidad desde el diseño.
  \item Diseñar e implementar el componente de detección, seguimiento y asociación de identidad
  entre cámaras, y documentar la granularidad de salida resultante frente al principio de
  proporcionalidad.
  \item Extender el prototipo con analítica de flujo peatonal (cruces de acceso, ocupación por
  zona, dashboard ejecutivo) sobre una cartografía real del terminal, evaluando su
  desempeño con pruebas automatizadas.
  \item Evaluar el desempeño del sistema mediante métricas cuantitativas y cualitativas,
  indicando en cada caso sobre qué conjunto de datos fue medido cada resultado.
\end{enumerate}

\clearpage
\section{Planteamiento del Problema y Análisis del Contexto}
\label{sec:planteamiento}

\subsection{Descripción del Caso}
\label{subsec:descripcion-caso}

El Aeropuerto Internacional Jorge Chávez no distribuye a sus pasajeros de manera
uniforme. En determinadas franjas horarias y en puntos concretos del terminal se forman
concentraciones muy superiores a las del resto del recinto en ese mismo momento. Medir esa
combinación (cuántas personas, en qué zona, durante cuánto tiempo) importa porque las
decisiones que la atienden son de corto plazo y reversibles: habilitar un mostrador adicional,
redirigir el flujo hacia un filtro descargado, adelantar un relevo de personal o abrir una segunda
fila. Todas dependen de conocer dónde se está formando la concentración mientras se forma; un
recuento que llega al cierre de la jornada sirve para planificar el mes siguiente, pero no para
resolver la cola de esa mañana.

\subsection{Identificación de Variables y Restricciones}
\label{subsec:variables-restricciones}

\subsubsection{Variables Operativas de Salida}

\begin{itemize}
  \item \textbf{Variable de salida principal:} concentración de personas por zona (cuántas,
  dónde, durante cuánto tiempo).
  \item \textbf{Granularidad de salida:} admite dos niveles con consecuencias normativas
  distintas: un mapa de densidad agregado por zona, que no requiere identificar
  individuos, o una localización individual por puntos, que produce un dato
  espacio-temporal atribuible a una persona si se persiste y se encadena en el tiempo.
\end{itemize}

\subsubsection{Variables Técnicas del Pipeline}

Parámetros reales configurados en \texttt{proyecto\_lap\_prototipo/}:

\begin{itemize}
  \item \textbf{Detección:} umbral de confianza \texttt{umbral\_deteccion = 0.5}; resolución de
  entrada del detector configurable en 480, 640 o 960 px.
  \item \textbf{Seguimiento mono-cámara:} asociación húngara de alta confianza con
  \texttt{DISTANCIA\_MAX\_ALTA\_PX = 45.0}; filtro de Kalman de velocidad constante.
  \item \textbf{Fusión entre cámaras:} \texttt{umbral\_distancia\_fusion = 1.5 m};
  \texttt{ventana\_temporal\_s = 1.0 s} de recencia.
  \item \textbf{Filtrado geométrico:} \texttt{margen\_px = 15.0} de tolerancia en el borde del
  polígono de alcance.
  \item \textbf{Descriptor de apariencia:} sustracción de fondo MOG2 (\texttt{history=200},
  \texttt{varThreshold=25}, \texttt{detectShadows=False}).
  \item \textbf{Monitoreo integrado:} muestreo de 0.2 s por análisis; hasta 32 cámaras
  configurables por servidor (sin garantía de procesamiento simultáneo de todas).
\end{itemize}

\subsubsection{Variables de Desempeño y Evaluación}

Medidas, no estimadas:

\begin{itemize}
  \item Tiempo de inferencia de P2PNet en CPU: $\approx$3.5 s por cada 0.2 s de video de
  origen.
  \item Filtrado de identificadores: de 29 identificadores brutos a 10 personas reales en la
  corrida de validación.
  \item Fusiones cruzadas verificadas entre cámara A y B: 2, con distancia espacio-temporal
  promedio 0.76 y máxima 1.18 sobre un umbral de 1.5.
  \item Suite de pruebas del monitoreo integrado: 65 pruebas automatizadas, más compilación
  de TypeScript y Vite.
\end{itemize}

\subsection{Restricciones del Proyecto}
\label{subsec:restricciones}

\begin{itemize}
  \item \textbf{Restricciones Técnicas.} Esta entrega corresponde a un desarrollo y validación
  conceptual sobre datos públicos y grabaciones propias, sin acceso todavía a material
  del Aeropuerto Jorge Chávez; la literatura revisada documenta una brecha de dominio
  conocida entre esos conjuntos y un entorno de despliegue nuevo. El sistema debe
  operar en CPU, sin GPU dedicada, sobre cámaras fijas que no comparten resolución,
  altura ni ángulo de instalación.
  \item \textbf{Restricciones Éticas.} El sistema no puede derivar del video ningún rasgo que
  permita identificar a una persona, ni siquiera de forma indirecta (RNF-01). Un
  identificador persistente que permite reagrupar apariciones de un mismo individuo es
  disociación, no anonimización, y el dato disociado sigue siendo dato personal.
  \item \textbf{Restricciones Legales.} La Ley 29733 y su reglamento (D.S. 016-2024-JUS)
  exigen privacidad desde el diseño y por defecto. La Directiva 01-2020-JUS/DGTAIPD
  fija plazos de conservación de imágenes de videovigilancia entre 30 y 60 días,
  mientras que el D.S. 007-2020-IN (D.L. 1218 y Ley 30120, aplicable a
  establecimientos con aforo de 50 personas o más, umbral que un terminal aéreo
  supera con holgura) exige un mínimo de 45 días; ambos regímenes son compatibles
  solo dentro de una ventana estrecha, porque responden a finalidades distintas.
\end{itemize}

\subsection{Análisis Técnico, Ético, Social, Ambiental}
\label{subsec:analisis-teas}

\begin{itemize}
  \item \textbf{Análisis técnico.} La arquitectura separa deliberadamente cuántas personas hay
  (disponible incluso sin calibrar cámaras) de quién es quién entre cámaras (que exige
  calibración y sincronización declarada). La auditoría del prototipo muestra, sin
  embargo, que la implementación actual sí persiste la información necesaria para
  reconstruir la segunda, lo que reabre la pregunta de si esa capacidad debe exponerse
  en producción.

  \item \textbf{Análisis ético y legal.} El artículo 2 de la Constitución Política del Perú (1993)
  reconoce el derecho a que los servicios informáticos no afecten la intimidad y el
  derecho a la imagen propia. El sistema trata cuatro tipos de dato con régimen
  distinto: el fotograma (dato personal), el punto con coordenadas (no identifica por sí
  solo), el dato espacio-temporal (que encadenado reconstruye una trayectoria) y el
  dato biométrico propiamente dicho, que ni un mapa de densidad ni un conjunto de
  coordenadas de cabezas constituyen por sí mismos. El equipo detectó una
  discrepancia terminológica relevante: describir como personas anonimizadas el
  resultado de un proceso que asigna un identificador persistente es impreciso, porque
  la anonimización debe ser irreversible. De las diez medidas de privacidad desde el
  diseño que adopta el proyecto (M-01 a M-10), al cierre de esta versión se cumple
  M-02 (no persistencia del fotograma) pero no M-03 (salida agregada por zona sin
  coordenadas individuales persistidas): es el único hallazgo de incumplimiento que el
  equipo puede cerrar por sí solo, sin depender de LAP.

  \item \textbf{Análisis social.} Un sistema que no toma decisiones sobre personas individuales
  puede producir un daño distributivo si subcuenta sistemáticamente a un grupo:
  estatura, vestimenta, equipaje voluminoso o sillas de ruedas son factores
  identificables, aunque no cuantificados por ningún trabajo revisado. Existe además
  un riesgo de sustitución progresiva del criterio humano por el del sistema, si un
  operador aprueba mecánicamente las alertas automáticas.

  \item \textbf{Análisis ambiental.} El equipo no ha medido el consumo energético del sistema.
  Reutilizar modelos preentrenados sin reentrenar, procesar en el borde sin transmitir
  el fotograma completo, y no persistir imágenes de forma permanente reducen, al
  mismo tiempo, la exposición de datos personales y el cómputo necesario.
\end{itemize}

\clearpage
\section{Marco Teórico}
\label{sec:marco-teorico}

El marco teórico revisa cuatro trabajos que sostienen, en conjunto, las decisiones técnicas
del prototipo: dos resuelven el problema de cuántos y dónde (conteo y localización), y dos
resuelven el problema de quién sigue siendo quién a lo largo del tiempo (seguimiento).

\subsection{CSRNet: redes convolucionales dilatadas para conteo de multitudes (Li, Zhang
y Chen, 2018)}
\label{subsec:csrnet}

\textbf{Qué propone.} Una red completamente convolucional con un \emph{front-end} que
reutiliza las primeras diez capas de VGG-16 y un \emph{back-end} de convoluciones dilatadas
que amplía el campo receptivo sin \emph{pooling}. La salida es un mapa de densidad continuo,
no un número ni una posición individual.

\textbf{Resultados.} Sobre ShanghaiTech, reduce el MAE en 7\,\% respecto de CP-CNN en la
Parte A y en 47.3\,\% en la Parte B; sobre TRANCOS, 15.4\,\% menos que el mejor método
previo.

\textbf{Por qué aplica al caso LAP.} Su salida es agregada por construcción, no requiere
rasgos faciales; es puramente convolucional y admite resoluciones de entrada variables,
relevante porque las cámaras de un terminal no comparten resolución, altura ni ángulo.

\textbf{Limitación reconocida por los autores.} En escenas dispersas (UCSD), CSRNet
queda por detrás de la mayoría de métodos salvo MCNN en MAE: su ventaja se concentra en
alta densidad, y buena parte de la jornada de un terminal aéreo transcurre en baja densidad.

\subsection{P2PNet: localización pura por puntos (Song et al., 2021)}
\label{subsec:p2pnet}

\textbf{Qué propone.} Un marco que predice directamente un conjunto de puntos con
coordenadas y confianza, sin representación intermedia, con emparejamiento uno a uno vía
algoritmo húngaro; aporta además la métrica \emph{density Normalized Average Precision}
(nAP).

\textbf{Resultados.} Sobre ShanghaiTech Parte A, reduce el MAE 4.8\,\% y el MSE 12.9\,\%
frente a ADSCNet; sobre UCF\_CC\_50 rebaja el MAE en 2.1 puntos; MAE de 85.32 sobre
UCF-QNRF, con nAP a 0.5 cercano al 90\,\% en la mayoría de los conjuntos.

\textbf{Por qué aplica al caso LAP.} La localización individual convierte una cifra en una
decisión operativa, pero eleva el perfil de riesgo normativo: un punto persistido y encadenado
en el tiempo se acerca a un seguimiento de trayectorias, una finalidad distinta de estimar
aglomeraciones.

\textbf{Limitación reconocida por los autores.} En UCF-QNRF su precisión no es tan
competitiva frente a ADSCNet; en NWPU-Crowd su desempeño es inferior por depender de un
único mapa de características de una sola escala.

\subsection{P2R: pérdida punto a región para conteo semi-supervisado (Lin, Zhao y Chan,
2025)}
\label{subsec:p2r}

\textbf{Qué propone.} Un marco semi-supervisado maestro-estudiante que sustituye el
emparejamiento punto a punto por uno punto a región, eliminando la necesidad del algoritmo
húngaro en esa etapa.

\textbf{Resultados.} Con solo 5\,\% de datos anotados, supera al aprendizaje totalmente
supervisado con 10\,\% (MAE 93.7, MSE 155.2 sobre ShanghaiTech Parte A); bajo supervisión
completa, supera a P2PNet en los cuatro conjuntos evaluados, con un cómputo de pérdida
$\approx$68 veces más rápido por imagen.

\textbf{Por qué aplica al caso LAP.} Conecta con el principio de minimización de datos de la
Ley 29733: un método competitivo con 5\,\% de datos etiquetados reduce en proporción similar
el volumen de imágenes que deberían anotarse manualmente.

\textbf{Limitación reconocida por los autores.} Fallo por desenfoque de movimiento de la
multitud, la norma y no la excepción en un terminal aéreo; los autores documentan
explícitamente la brecha de dominio.

\subsection{ByteTrack: seguimiento multiobjeto asociando toda caja de detección (Zhang
et al., 2022)}
\label{subsec:bytetrack}

\textbf{Qué propone.} Asociar todas las cajas de detección, no solo las de alta confianza:
para las de baja puntuación, aprovecha su similitud con las trayectorias activas para recuperar
objetos reales y descartar detecciones de fondo, combinado con un filtro de Kalman.

\textbf{Resultados.} Desempeño de referencia en los benchmarks MOT20, HiEve y
BDD100K.

\textbf{Por qué aplica al caso LAP.} El prototipo adapta esta lógica a puntos en el
seguimiento mono-cámara (Kalman de velocidad constante más asignación húngara en dos
etapas), y la extensión de monitoreo integrado usa ByteTrack sin modificar sobre cajas de un
detector YOLO11n.

\subsection{Síntesis de antecedentes}
\label{subsec:sintesis-antecedentes}

CSRNet resuelve cuántos con una superficie de densidad de la que no puede extraerse
dónde. P2PNet resuelve dónde eliminando la representación intermedia, pero paga ese avance
con un costo de anotación creciente. P2R resuelve ese costo llevando el enfoque a régimen
semi-supervisado. ByteTrack resuelve quién sigue siendo quién entre cuadros consecutivos sin
descartar evidencia débil. De la lectura conjunta se derivan tres criterios que gobiernan el
diseño del sistema: la granularidad de la salida es una decisión de cumplimiento normativo y
no solo de producto; la eficiencia de etiquetado de P2R es, además de una ventaja técnica, un
argumento de minimización de datos; y ninguna cifra de la literatura revisada es transferible al
Aeropuerto Jorge Chávez sin una validación local que hoy no se ha realizado.

\clearpage
% Ficha de caso de uso:
%   \fichacu{ID}{Nombre}{Actor}{Precondicion}{items del flujo}{Postcondicion}{Excepcion o nota}
\newcommand{\fichacu}[7]{%
  \subsubsection*{#1. #2}
  \noindent\textbf{Actor:} #3\par
  \noindent\textbf{Precondición:} #4\par
  \noindent\textbf{Flujo principal:}
  \begin{enumerate}[leftmargin=2.2em,itemsep=0.1em,topsep=0.2em]
  #5
  \end{enumerate}
  \noindent\textbf{Postcondición:} #6\par
  \if\relax\detokenize{#7}\relax\else\noindent\textbf{Excepción o nota:} #7\par\fi
  \medskip}

\section{Especificación de Requerimientos}
\label{sec:requerimientos}

Este apartado describe el prototipo tal como está implementado, no una versión aspiracional.
Adopta la numeración de la hoja maestra de requerimientos del equipo: dieciocho casos de uso
(CU-01 a CU-18), dieciocho requisitos funcionales (RF-01 a RF-18) y ocho requisitos no
funcionales (RNF-01 a RNF-08), relacionados en la matriz de trazabilidad del
apartado~\ref{subsec:trazabilidad}. El estado de cada requisito se contrastó con el código del
prototipo vigente. Donde la hoja maestra describía una versión anterior del sistema (el detector
seleccionable en CU-03 y RF-03, el desempate por color en CU-04, la cifra de inferencia de
RNF-06 y el estado de RNF-01), este informe recoge el comportamiento actual. La
reidentificación por apariencia y la asociación entre cámaras se documentan en CU-04 y RF-04.

\subsection{Requerimientos funcionales}
\label{subsec:rf}

\subsubsection{Casos de uso del sistema}
\label{subsec:casos-uso-tabla}

La Tabla~\ref{tab:casos-uso} resume los dieciocho casos de uso y la
Figura~\ref{fig:casos-uso} los relaciona con sus actores. El caso CU-16 conserva su número en
la hoja maestra, pero corresponde a un módulo retirado del prototipo vigente.

{\small
\setlength{\LTleft}{0pt}\setlength{\LTright}{0pt}
\begin{longtable}{@{}L{1.3cm}L{6.0cm}L{3.8cm}L{2.6cm}@{}}
\caption{Casos de uso del sistema AeroTrack.}\label{tab:casos-uso}\\
\toprule
\textbf{ID} & \textbf{Caso de uso} & \textbf{Actor} & \textbf{Estado} \\
\midrule
\endfirsthead
\toprule
\textbf{ID} & \textbf{Caso de uso} & \textbf{Actor} & \textbf{Estado} \\
\midrule
\endhead
\bottomrule
\endlastfoot
CU-01 & Configurar el plano y las cámaras & Operador & Implementado \\
CU-02 & Calibrar una cámara & Operador & Implementado \\
CU-03 & Procesar una fuente y hacer seguimiento & Operador & Implementado \\
CU-04 & Asociar identidad entre cámaras y reidentificar personas & Sistema & Implementado \\
CU-05 & Detectar y alertar aglomeraciones (regla global) & Operador o administrador & Implementado \\
CU-06 & Visualizar el plano en vivo y el video por cámara & Operador o administrador & Implementado \\
CU-07 & Exportar reportes & Operador o administrador & Implementado \\
CU-08 & Gestionar zonas de interés & Operador & Implementado \\
CU-09 & Configurar reglas de aglomeración por zona & Operador & Implementado \\
CU-10 & Consultar el dashboard ejecutivo y el panel comercial & Operador o administrador & Implementado \\
CU-11 & Validar una fuente en modo laboratorio & Operador & Implementado \\
CU-12 & Consultar la auditoría del sistema & Operador o administrador & Implementado \\
CU-13 & Gestionar proyectos independientes & Operador o administrador & Implementado \\
CU-14 & Iniciar sesión según el rol & Operador o administrador & Implementado \\
CU-15 & Revisar la cámara en vivo antes de calibrar & Operador & Implementado \\
CU-16 & Configurar y recibir avisos de equipaje sin custodia & No aplica & Retirado \\
CU-17 & Configurar y recibir alertas de seguridad por correo & Operador o administrador; personal de seguridad & Implementado \\
CU-18 & Gestionar la bitácora de incidentes & Operador o administrador & Implementado \\
\end{longtable}
}

\begin{figure}[H]
\centering
\resizebox{\anchofigura}{!}{%
\begin{tikzpicture}[
  cu/.style={ellipse,draw=black,thick,minimum width=5.6cm,minimum height=1.05cm,
             text width=4.5cm,align=center,font=\footnotesize,inner sep=0pt},
  retirado/.style={cu,dashed},
  asoc/.style={black,thin}]
  \draw[thick] (-1.0,0.9) rectangle (13.6,-15.58);
  \node[font=\small\bfseries,anchor=north] at (6.3,0.8) {Sistema AeroTrack};
  \node[cu] (c01) at (2.0,-1.42) {CU-01\\Configurar plano y cámaras};
  \node[cu] (c02) at (2.0,-3.04) {CU-02\\Calibrar una cámara};
  \node[cu] (c03) at (2.0,-4.66) {CU-03\\Procesar fuente y seguir};
  \node[cu] (c08) at (2.0,-6.28) {CU-08\\Gestionar zonas de interés};
  \node[cu] (c09) at (2.0,-7.90) {CU-09\\Reglas por zona};
  \node[cu] (c11) at (2.0,-9.52) {CU-11\\Validar en laboratorio};
  \node[cu] (c15) at (2.0,-11.14) {CU-15\\Revisar cámara en vivo};
  \node[cu] (c05) at (10.6,-1.42) {CU-05\\Alertar aglomeraciones};
  \node[cu] (c06) at (10.6,-3.04) {CU-06\\Visualizar plano y video};
  \node[cu] (c07) at (10.6,-4.66) {CU-07\\Exportar reportes};
  \node[cu] (c10) at (10.6,-6.28) {CU-10\\Dashboard y panel comercial};
  \node[cu] (c12) at (10.6,-7.90) {CU-12\\Consultar auditoría};
  \node[cu] (c13) at (10.6,-9.52) {CU-13\\Gestionar proyectos};
  \node[cu] (c14) at (10.6,-11.14) {CU-14\\Iniciar sesión por rol};
  \node[cu] (c17) at (10.6,-12.76) {CU-17\\Alertas por correo};
  \node[cu] (c18) at (10.6,-14.38) {CU-18\\Bitácora de incidentes};
  \node[cu] (c04) at (2.0,-12.76) {CU-04\\Identidad y reidentificación};
  \node[retirado] (c16) at (4.6,-14.38) {CU-16 (retirado)\\Equipaje sin custodia};
  \begin{scope}[shift={(-2.90,-6.48)}]
    \draw[thick] (0,0.55) circle (0.22); \draw[thick] (0,0.33)--(0,-0.35);
    \draw[thick] (-0.35,0.1)--(0.35,0.1); \draw[thick] (0,-0.35)--(-0.3,-0.85); \draw[thick] (0,-0.35)--(0.3,-0.85);
    \node[font=\small,below,align=center] at (0,-0.9) {Operador};
  \end{scope}
  \begin{scope}[shift={(15.60,-8.10)}]
    \draw[thick] (0,0.55) circle (0.22); \draw[thick] (0,0.33)--(0,-0.35);
    \draw[thick] (-0.35,0.1)--(0.35,0.1); \draw[thick] (0,-0.35)--(-0.3,-0.85); \draw[thick] (0,-0.35)--(0.3,-0.85);
    \node[font=\small,below,align=center] at (0,-0.9) {Operador o\\administrador};
  \end{scope}
  \begin{scope}[shift={(-2.90,-13.36)}]
    \draw[thick] (0,0.55) circle (0.22); \draw[thick] (0,0.33)--(0,-0.35);
    \draw[thick] (-0.35,0.1)--(0.35,0.1); \draw[thick] (0,-0.35)--(-0.3,-0.85); \draw[thick] (0,-0.35)--(0.3,-0.85);
    \node[font=\small,below,align=center] at (0,-0.9) {Sistema\\(automático)};
  \end{scope}
  \coordinate (op) at (-2.5,-6.38);
  \coordinate (ad) at (15.2,-8.00);
  \coordinate (si) at (-2.5,-13.26);
  \foreach \id in {01,02,03,08,09,11,15} { \draw[asoc] (op) -- (c\id.west); }
  \foreach \id in {05,06,07,10,12,13,14,17,18} { \draw[asoc] (ad) -- (c\id.east); }
  \draw[asoc] (si) -- (c04.west);
  \draw[asoc,dashed] (si) -- (c16.west);
\end{tikzpicture}}
\caption{Diagrama de casos de uso del sistema AeroTrack. La línea discontinua señala el
módulo retirado del prototipo vigente.}
\label{fig:casos-uso}
\end{figure}

\subsubsection{Fichas de casos de uso}
\label{subsec:casos-uso-detalle}

\fichacu{CU-01}{Configurar el plano y las cámaras}
  {Operador.}
  {el servidor local está en ejecución, el usuario inició sesión con el rol de operador y hay un
   proyecto abierto.}
  {\item El operador elige uno de los planos del aeropuerto, crea un plano en blanco o importa
         una imagen, un PDF o un archivo DXF.
   \item Si importa una imagen o un PDF, recorta la región útil y el sistema la reduce a un dibujo
         esquemático de líneas por detección de bordes, en lugar de conservar la fotografía.
   \item Coloca cada cámara sobre el plano y define su alcance orientativo.
   \item Declara qué cámaras están relacionadas entre sí.}
  {el plano queda registrado con sus cámaras y disponible para la calibración (CU-02).}
  {}

\fichacu{CU-02}{Calibrar una cámara}
  {Operador.}
  {la cámara tiene una fuente de video válida.}
  {\item El operador marca un punto de referencia del suelo en el video y su correspondencia en
         el plano.
   \item Repite la marcación un mínimo de cuatro veces; se recomiendan de seis a ocho puntos
         repartidos por toda el área de interés.
   \item El sistema calcula la homografía y, con cinco o más puntos, comprueba cada referencia
         dejándola fuera del ajuste.}
  {la cámara queda habilitada para proyectar posiciones al plano compartido.}
  {el sistema rechaza la calibración si hay referencias repetidas, casi alineadas o
   geométricamente inconsistentes. La calibración es opcional: sin ella la cámara detecta y sigue
   personas en la imagen, pero no las ubica en el plano ni usa la posición para asociarlas entre
   cámaras.}

\fichacu{CU-03}{Procesar una fuente y hacer seguimiento}
  {Operador.}
  {la cámara tiene una fuente configurada (archivo de video, cámara USB o flujo de red) y su zona
   útil marcada; los pesos del detector YOLO y el modelo OSNet están instalados en el equipo.}
  {\item El operador inicia la sesión de monitoreo y elige el modo de rendimiento (automático,
         preciso, equilibrado o rápido).
   \item El sistema elige la variante de YOLO11 según el hardware del equipo y detecta personas en
         cada muestra de video.
   \item Descarta las detecciones fuera de la zona útil, las dudosas mucho más pequeñas que las
         personas seguras de la misma cámara y los objetos inmóviles de baja confianza.
   \item Sigue a cada persona con asociación en dos etapas sobre el punto de los pies, reforzada
         por la superposición de la caja predicha y una firma de color del torso, y le asigna un
         identificador local estable dentro de esa cámara.
   \item Si la cámara está calibrada, proyecta el punto de los pies al plano.}
  {cada persona detectada tiene un identificador local en esa cámara y, si la cámara está
   calibrada, una posición en el plano; las vistas confiables pasan al motor de identidad
   (CU-04).}
  {si falta el modelo OSNet, la sesión no se inicia. Si faltan los pesos de la variante de YOLO
   recomendada, el sistema usa la mejor variante disponible en el equipo y lo advierte.}

\fichacu{CU-04}{Asociar identidad entre cámaras y reidentificar personas}
  {Sistema (automático), a partir de la declaración de sincronización que hace el operador.}
  {sesión de monitoreo en curso. Para asociar personas entre cámaras: al menos dos cámaras
   relacionadas y la declaración expresa del operador de que los relojes de las fuentes están
   sincronizados. La calibración es opcional.}
  {\item Para cada persona seguida cuya vista es confiable (buena confianza, altura suficiente,
         lejos del borde y sin oclusión), el sistema calcula un vector de apariencia OSNet de 512
         dimensiones sobre el recorte de cuerpo entero; las partes tapadas por otra persona se
         rellenan con un gris neutro.
   \item Agrupa las vistas de cada tramo continuo de seguimiento y forma con ellas un prototipo
         de apariencia.
   \item La persona lleva un identificador provisional hasta reunir al menos cinco vistas durante
         dos segundos; entonces recibe un identificador confirmado, que es el único que cuenta en
         ocupación, flujo y alertas.
   \item Busca, entre las identidades de las cámaras relacionadas, una candidata con parecido de
         apariencia sobre el umbral y coherente en tiempo y, si hay calibración, en posición en el
         plano. Si hay una candidata inequívoca, reutiliza su identificador; si la mejor candidata es
         ambigua, marca la asociación como incierta en lugar de forzarla.
   \item Si la memoria de apariencia está activa, que es la configuración por defecto, compara
         cada identidad nueva con todas las personas guardadas dentro del plazo de retención y,
         si coincide sin ambigüedad, le devuelve su identificador anterior, aunque se trate de otra
         sesión u otro día.
   \item Al cerrar la sesión, reagrupa los tramos con la grabación completa a la vista y renumera
         a las personas contadas de 1 a N.}
  {cada observación queda con un identificador, provisional o confirmado, y un estado de
   asociación: local, estimada (une cámaras), reidentificada (recupera un tramo o una persona de
   la memoria) o incierta.}
  {sin sincronización declarada, el sistema no asocia personas entre cámaras en la sesión, pero
   la memoria de apariencia sigue comparando cada identidad nueva con las personas guardadas.
   La memoria conserva el vector en el disco del equipo durante 24 horas por defecto
   (configurable hasta 168), lo que constituye la reidentificación individual persistente que
   RNF-01 prohíbe (apartado~\ref{subsec:privacidad}).}

\fichacu{CU-05}{Detectar y alertar aglomeraciones (regla global)}
  {Operador o administrador.}
  {sesión de monitoreo en ejecución.}
  {\item Desde la vista de alertas, el usuario fija un radio en metros, una cantidad mínima de
         personas y un tiempo de permanencia, aplicables a todo el plano.
   \item El sistema evalúa en cada muestra las posiciones de las personas confirmadas.
   \item Marca una alerta cuando una concentración cumple los tres criterios de forma
         sostenida.}
  {la alerta queda registrada como incidente en la bitácora (CU-18) y se envía por correo si el
   aviso está configurado (CU-17).}
  {}

\fichacu{CU-06}{Visualizar el plano en vivo y el video por cámara}
  {Operador o administrador.}
  {sesión activa o grabación de una sesión anterior disponible.}
  {\item El usuario observa el plano con las personas, sus trayectorias recientes y las
         aglomeraciones resaltadas.
   \item Alterna al video real de una cámara, con la caja y el identificador de cada persona
         superpuestos.
   \item Puede revisar las sesiones grabadas desde la vista de videos y resultados.}
  {la consulta no modifica la configuración ni la sesión.}
  {}

\fichacu{CU-07}{Exportar reportes}
  {Operador o administrador.}
  {sesión con datos observados.}
  {\item El usuario elige el tipo de reporte: zonas de ocupación, comparación de sectores,
         ocupación por zona, flujo horario, trayectorias recientes o reporte comercial.
   \item Elige el formato: PDF, Excel, CSV o JSON.}
  {el sistema genera el archivo indicando siempre el origen de los datos: simulación sintética,
   videos de prueba o fuente real sin validación de precisión.}
  {el reporte de trayectorias contiene recorridos por identificador, por lo que hereda la
   calificación de dato personal de las grabaciones (Anexo~\ref{anx:marco}).}

\fichacu{CU-08}{Gestionar zonas de interés}
  {Operador.}
  {plano configurado.}
  {\item El operador dibuja polígonos, rectángulos o círculos sobre el plano.
   \item Clasifica cada área como zona de interés, cola, zona restringida, sala, pasillo, muro,
         puerta o zona comercial.
   \item Puede aceptar puntos sugeridos por el sistema según la ocupación observada.}
  {las zonas quedan registradas y disponibles para reglas y reportes por zona.}
  {}

\fichacu{CU-09}{Configurar reglas de aglomeración por zona}
  {Operador.}
  {al menos una zona dibujada.}
  {\item El operador activa una regla independiente para una zona, con su propia cantidad mínima
         de personas y su propio tiempo de permanencia.}
  {la zona genera alertas específicas según su uso; por ejemplo, una cola tolera más personas
   que una sala de espera.}
  {}

\fichacu{CU-10}{Consultar el dashboard ejecutivo y el panel comercial}
  {Operador o administrador.}
  {sesión con datos observados.}
  {\item El usuario revisa el ranking de zonas por ocupación acumulada, la evolución de las
         personas observadas y las zonas sugeridas para evaluación comercial.
   \item En el panel comercial consulta las zonas ordenadas por oportunidad comercial, el
         tráfico por acceso y las métricas de cada negocio (visitas, permanencias y regresos).}
  {la información se presenta agregada y el panel indica si procede de una fuente real o es
   sintética.}
  {aunque se presentan agregadas, las métricas por negocio se calculan a partir de las
   trayectorias por identificador, lo que constituye una finalidad distinta de la estimación de
   aglomeraciones (Sección~\ref{subsec:impacto-etico}).}

\fichacu{CU-11}{Validar una fuente en modo laboratorio}
  {Operador.}
  {cámara agregada al proyecto.}
  {\item El operador prueba la cámara de forma aislada: señal, resolución, frecuencia de cuadros,
         tiempo de procesamiento, calidad de la calibración y estabilidad del seguimiento, con
         identificadores temporales sobre el video.}
  {la cámara queda validada o rechazada antes de sumarse a una sesión real.}
  {}

\fichacu{CU-12}{Consultar la auditoría del sistema}
  {Operador o administrador.}
  {al menos una acción registrada.}
  {\item El usuario revisa el registro de acciones de la ejecución del servidor (guardar, iniciar,
         pausar, detener, abrir una vista en vivo o borrar la memoria de identidad, entre otras),
         con fecha, hora y detalle.}
  {queda disponible la trazabilidad operativa de la sesión, sin identidades reales.}
  {}

\fichacu{CU-13}{Gestionar proyectos independientes}
  {Operador; el administrador solo puede abrir proyectos existentes.}
  {servidor en ejecución con al menos un proyecto.}
  {\item El usuario ve la lista de proyectos desde la pantalla de inicio o el selector de la
         cabecera.
   \item El operador crea un proyecto nuevo, abre uno existente, lo renombra o lo elimina.}
  {el sistema cambia el proyecto activo; el plano, las cámaras, la calibración, los incidentes, las
   ventas y la memoria de apariencia de un proyecto nunca aparecen en otro.}
  {no se puede eliminar el único proyecto existente, y el sistema rechaza los cambios destinados
   a un proyecto distinto del abierto.}

\fichacu{CU-14}{Iniciar sesión según el rol}
  {Operador o administrador.}
  {existe al menos un usuario; si no existe ninguno, el primero que se registra queda como
   operador.}
  {\item El usuario ingresa su nombre y su contraseña.
   \item El sistema valida las credenciales contra el resumen criptográfico almacenado y abre una
         sesión con el rol correspondiente.}
  {el operador accede a la configuración y a todo lo demás; el administrador accede a todo
   menos a la configuración y puede ajustar los umbrales de alerta desde la vista de alertas.}
  {las credenciales incorrectas se rechazan sin indicar cuál de los dos datos falló.}

\fichacu{CU-15}{Revisar la cámara en vivo antes de calibrar}
  {Operador.}
  {la cámara tiene una fuente configurada y no hay una sesión de seguimiento en curso.}
  {\item El operador abre la vista en vivo de la cámara desde la etapa de cámaras del asistente.
   \item El sistema muestra la imagen en movimiento sin ejecutar detección; si la fuente es una
         grabación, permite pausar, reanudar y saltar a un instante.
   \item El operador puede obtener una imagen fija para marcar la zona útil.}
  {el operador confirma encuadre, luz y retardo; la apertura queda en el registro de auditoría y,
   con la imagen fija, la fuente queda comprobada con su resolución y su frecuencia de cuadros.}
  {si nadie solicita imágenes durante 20 segundos, el sistema libera la cámara.}

\fichacu{CU-16}{Configurar y recibir avisos de equipaje sin custodia}
  {no aplica en el prototipo vigente.}
  {no aplica.}
  {\item Este caso describía la detección de mochilas, bolsos y maletas inmóviles durante un
         tiempo configurable y la emisión de un aviso para revisión humana.}
  {no aplica.}
  {el módulo se retiró del prototipo al simplificarlo al esquema de seguimiento multicámara. Al
   abrir proyectos antiguos, su configuración se elimina. Se conserva el número para no alterar
   la numeración maestra. Los incidentes y los correos actuales se refieren solo a
   aglomeraciones.}

\fichacu{CU-17}{Configurar y recibir alertas de seguridad por correo}
  {Operador o administrador, que configura; personal de seguridad, que recibe.}
  {se cuenta con una cuenta de correo remitente y su contraseña de aplicación.}
  {\item Desde la vista de alertas, el usuario activa el aviso por correo e indica la cuenta
         remitente, su contraseña de aplicación y hasta veinte destinatarios.
   \item Puede enviar un correo de prueba.
   \item Ante cada incidente nuevo de aglomeración, el sistema envía un correo con el lugar, la
         cámara, el motivo del aviso y la advertencia de que requiere comprobación humana.}
  {la configuración queda guardada fuera de los archivos del proyecto, la contraseña no se
   muestra nunca en la interfaz y el cambio queda en el registro de auditoría.}
  {una misma zona no genera más de un correo cada 120 segundos y un aviso ya enviado no se
   reenvía; un fallo de envío queda registrado sin interrumpir la sesión.}

\fichacu{CU-18}{Gestionar la bitácora de incidentes}
  {Operador o administrador.}
  {existen incidentes de aglomeración registrados.}
  {\item El usuario consulta la lista de incidentes del proyecto activo, con su hora, zona,
         magnitud y estado.
   \item Marca cada incidente como revisado, falsa alarma o resuelto, o lo devuelve a
         pendiente.}
  {el incidente queda con su estado actualizado, lo que distingue las alertas atendidas de las
   pendientes.}
  {los incidentes se conservan aunque la sesión termine, y los de un proyecto nunca aparecen
   en otro.}

\subsubsection{Requisitos funcionales}
\label{subsec:rf-tabla}

{\small
\setlength{\LTleft}{0pt}\setlength{\LTright}{0pt}
\begin{longtable}{@{}L{1.3cm}L{6.6cm}L{4.2cm}L{1.9cm}@{}}
\caption{Requisitos funcionales y su estado verificado en el prototipo.}\label{tab:rf}\\
\toprule
\textbf{ID} & \textbf{Requisito} & \textbf{Estado} & \textbf{CU} \\
\midrule
\endfirsthead
\toprule
\textbf{ID} & \textbf{Requisito} & \textbf{Estado} & \textbf{CU} \\
\midrule
\endhead
\bottomrule
\endlastfoot
RF-01 & Configurar plano y cámaras, con recorte e importación como dibujo de líneas. & Implementado. & CU-01 \\
RF-02 & Calibrar por homografía con validación de consistencia geométrica. & Implementado. & CU-02 \\
RF-03 & Detectar personas con YOLO11, en la variante que corresponda al hardware, y seguirlas dentro de cada cámara con asociación en dos etapas sobre el punto de los pies. & Implementado; los pesos se instalan al preparar el sistema. & CU-03, CU-11 \\
RF-04 & Asociar identidad entre cámaras con el vector de apariencia OSNet y coherencia de tiempo y plano, solo entre cámaras relacionadas y con sincronización declarada, marcando \enquote{incierta} cuando corresponda, y reidentificar personas mediante la memoria de apariencia. & Implementado; su configuración por defecto incumple RNF-01. & CU-04 \\
RF-05 & Reglas de aglomeración globales configurables por radio, cantidad y permanencia. & Implementado. & CU-05 \\
RF-06 & Visualizar plano en vivo y video real por cámara con el identificador superpuesto. & Implementado. & CU-06 \\
RF-07 & Exportar reportes en PDF, Excel, CSV y JSON, por tipo (zonas, comparación de sectores, ocupación, flujo, trayectorias y comercial). & Implementado. & CU-07 \\
RF-08 & Dibujar y clasificar zonas de interés, incluyendo sugerencias automáticas por ocupación observada. & Implementado. & CU-08 \\
RF-09 & Configurar reglas de aglomeración independientes por zona. & Implementado. & CU-09 \\
RF-10 & Mostrar el dashboard ejecutivo con ranking de zonas y evolución temporal de personas observadas, y el panel comercial. & Implementado. & CU-10 \\
RF-11 & Validar una fuente de video en modo laboratorio antes de integrarla a la operación. & Implementado. & CU-11 \\
RF-12 & Registrar la auditoría de acciones del operador. & Implementado. & CU-12 \\
RF-13 & Gestionar proyectos independientes (crear, abrir, renombrar, eliminar), con aislamiento de datos entre ellos. & Implementado. & CU-13 \\
RF-14 & Autenticar usuarios con dos roles (operador y administrador) aplicados en el servidor, no solo ocultos en la interfaz. & Implementado. & CU-14 \\
RF-15 & Mostrar la vista en vivo de una cámara sin ejecutar detección, con controles de reproducción para grabaciones. & Implementado. & CU-15 \\
RF-16 & Detectar equipaje sin custodia por cámara, con tiempo de espera e intervalo de revisión configurables. & Retirado del prototipo vigente. & CU-16 \\
RF-17 & Enviar alertas de seguridad por correo electrónico, con control de repetición por incidente. & Implementado. & CU-17 \\
RF-18 & Registrar incidentes en una bitácora persistente con estado revisable (pendiente, revisado, falsa alarma, resuelto). & Implementado. & CU-18 \\
\end{longtable}
}

\subsection{Requerimientos de privacidad}
\label{subsec:privacidad}

La normativa de protección de datos exige que las garantías se incorporen desde la
concepción del tratamiento y que la configuración más protectora venga activada por defecto
(Sección~\ref{subsec:restricciones-legales}). El sistema recoge esa exigencia como requisitos no
funcionales, es decir, como condiciones que debe cumplir con independencia de su
funcionalidad.

El requisito rector es RNF-01, cuyo texto en la hoja maestra del equipo es: \enquote{No derivar
del video ningún rasgo que permita identificación biométrica o reidentificación individual
persistente}. \textbf{Al cierre de esta versión, RNF-01 es no conforme.} El sistema no construye
plantillas faciales, pero calcula para cada persona un vector de apariencia OSNet de cuerpo
entero y, con la memoria de apariencia activa, que es la configuración por defecto, lo guarda en
una base de datos local del equipo durante 24 horas (el plazo admite valores de 1 a 168 horas) y
lo usa para reconocer a la misma persona entre cámaras, entre sesiones y entre días. Eso es una
reidentificación individual persistente. La corrección está identificada y depende solo del
equipo: desactivar por defecto la memoria de apariencia, de modo que el vector viva
únicamente en memoria durante la sesión y se descarte con ella. Queda como acción pendiente.

De RNF-01 y del principio de proporcionalidad se derivan diez medidas de privacidad desde el
diseño. La Tabla~\ref{tab:medidas} las recoge con su estado en el prototipo.

{\small
\setlength{\LTleft}{0pt}\setlength{\LTright}{0pt}
\begin{longtable}{@{}L{1.2cm}L{5.8cm}L{7.4cm}@{}}
\caption{Medidas de privacidad desde el diseño: lo exigido y lo implementado.}\label{tab:medidas}\\
\toprule
\textbf{ID} & \textbf{Medida exigida} & \textbf{Estado en el prototipo} \\
\midrule
\endfirsthead
\toprule
\textbf{ID} & \textbf{Medida exigida} & \textbf{Estado en el prototipo} \\
\midrule
\endhead
\bottomrule
\endlastfoot
M-01 & Procesamiento en el borde, de modo que el fotograma no salga del perímetro del terminal. & Cumplida en el prototipo: todo el procesamiento ocurre en el equipo local y el servidor solo acepta conexiones del propio equipo. La localización del despliegue final sigue pendiente con LAP. \\
M-02 & No persistencia del fotograma, que se procesa en memoria y se descarta. & Cumplida para las cámaras en vivo. Los videos cargados para pruebas se conservan en el equipo. \\
M-03 & Salida agregada por zona, sin coordenadas individuales persistidas. & No cumplida: las grabaciones de cada sesión conservan, por muestra y por persona, el identificador, la caja y la posición en el plano, sin plazo de retención automático. \\
M-04 & Difuminado irreversible de rostros si algún fotograma debe conservarse. & No implementada: los videos de prueba se conservan sin difuminar. \\
M-05 & Plazo de retención definido, con purga automática. & Parcial: la memoria de apariencia se purga sola tras 24 horas; las grabaciones de sesión no tienen plazo automático y se borran a mano. \\
M-06 & Control de acceso por roles, con registro de auditoría. & Cumplida. \\
M-07 & Cifrado en tránsito y en reposo. & No implementada: la conexión con la consola es local y sin cifrar, y las bases de datos locales no están cifradas. \\
M-08 & Carteles informativos conforme al Anexo~\ref{anx:cartel}. & Pendiente: el modelo está propuesto, pero sus campos dependen de LAP. \\
M-09 & Documentación del linaje de los datos. & Parcial: el informe describe qué se calcula y dónde se guarda, sin un registro formal de linaje. \\
M-10 & Supervisión humana: el sistema alerta y no ejecuta acciones automáticas sobre personas. & Cumplida. \\
\end{longtable}
}

Donde un parámetro admita un ajuste más protector y otro menos protector (resolución de
captura, frecuencia de muestreo, granularidad de la salida, memoria de apariencia, plazo de
retención), el valor predeterminado debe ser el más protector. \textbf{Estas medidas describen el
sistema exigido y no el implementado}: una lista de medidas presentada al cliente como estado
alcanzado, cuando es estado deseado, es exactamente el tipo de afirmación que este informe
evita. A ello se suma una carencia de interfaz: el borrado inmediato de la memoria de identidad
existe como servicio interno del servidor, pero la consola no ofrece ningún botón para
ejecutarlo, de modo que un operador no puede atender una solicitud de supresión sin
asistencia técnica.

\subsection{Requerimientos no funcionales}
\label{subsec:rnf}

{\small
\setlength{\LTleft}{0pt}\setlength{\LTright}{0pt}
\begin{longtable}{@{}L{1.4cm}L{5.0cm}L{5.8cm}L{1.9cm}@{}}
\caption{Requisitos no funcionales y su estado verificado en el prototipo.}\label{tab:rnf}\\
\toprule
\textbf{ID} & \textbf{Requisito} & \textbf{Estado} & \textbf{CU} \\
\midrule
\endfirsthead
\toprule
\textbf{ID} & \textbf{Requisito} & \textbf{Estado} & \textbf{CU} \\
\midrule
\endhead
\bottomrule
\endlastfoot
RNF-01 & No derivar del video ningún rasgo que permita identificación biométrica o reidentificación individual persistente. & \textbf{No conforme.} La memoria de apariencia, activa por defecto, guarda el vector OSNet de cuerpo entero 24 horas (hasta 168) y reidentifica entre sesiones y días; además, las grabaciones persisten posiciones individuales (M-03). Acción pendiente: desactivar la memoria por defecto y fijar una retención automática para las grabaciones. & CU-03, CU-04, CU-06, CU-07, CU-10 \\
RNF-02 & Servidor accesible solo desde el propio equipo, con credencial de sesión en toda operación de escritura. & Implementado. & Transversal \\
RNF-03 & No persistir imágenes de forma permanente. & Cumplido para fuentes en vivo: los fotogramas se procesan en memoria y se descartan. Los videos cargados para pruebas se conservan en el equipo y su retención queda pendiente de definir. & CU-03, CU-06 \\
RNF-04 & Declarar expresamente cuándo el sistema opera en modo de demostración sintética. & Implementado. & CU-06, CU-07, CU-10 \\
RNF-05 & No presentar ninguna cifra de desempeño sobre datos públicos ni sobre grabaciones propias como desempeño esperado en el Aeropuerto Jorge Chávez sin validación en sitio. & Aplicado en este informe (Secciones~\ref{subsec:variables-restricciones} y~\ref{subsec:analisis-teas}). & CU-03, CU-11 \\
RNF-06 & Tiempo de inferencia compatible con la operación práctica. & No verificado como cumplido. Con YOLO11 nano y OSNet en CPU, analizar varias cámaras con una muestra cada 0,2 segundos va, según estimación del equipo, entre cinco y veinticinco veces más lento que el video; los modos de rendimiento espacian las muestras. Falta una medición con protocolo documentado. & CU-03, CU-11 \\
RNF-07 & Almacenar las contraseñas con resumen criptográfico y sal por usuario, nunca en texto plano ni expuestas por la interfaz de programación. & Implementado (PBKDF2 con SHA-256 y 240\,000 iteraciones). & CU-14 \\
RNF-08 & Los datos de un proyecto (plano, cámaras, incidentes, ventas, sesiones grabadas y memoria de apariencia) no deben ser accesibles desde otro proyecto. & Implementado y verificado con pruebas automatizadas. & CU-13, CU-18 \\
\end{longtable}
}

\subsection{Matriz de trazabilidad}
\label{subsec:trazabilidad}

La Tabla~\ref{tab:trazabilidad} relaciona cada requisito con los casos de uso que lo realizan.
Reproduce la matriz de la hoja maestra con una corrección: RNF-01 se vincula a los casos de uso
que calculan, conservan o explotan identificadores y trayectorias por persona (CU-03, CU-04,
CU-06, CU-07 y CU-10), en lugar de a la gestión de zonas y reglas.

\begin{table}[H]
\centering
\footnotesize
\setlength{\tabcolsep}{2.2pt}
\caption{Matriz de trazabilidad entre requisitos y casos de uso (columnas CU-01 a CU-18).}
\label{tab:trazabilidad}
\begin{tabular}{@{}lC{0.42cm}C{0.42cm}C{0.42cm}C{0.42cm}C{0.42cm}C{0.42cm}C{0.42cm}C{0.42cm}C{0.42cm}C{0.42cm}C{0.42cm}C{0.42cm}C{0.42cm}C{0.42cm}C{0.42cm}C{0.42cm}C{0.42cm}C{0.42cm}@{}}
\toprule
 & \multicolumn{18}{c}{\textbf{Caso de uso (CU-)}} \\
\cmidrule(l){2-19}
\textbf{Req.} & 01 & 02 & 03 & 04 & 05 & 06 & 07 & 08 & 09 & 10 & 11 & 12 & 13 & 14 & 15 & 16 & 17 & 18 \\
\midrule
RF-01 & X &  &  &  &  &  &  &  &  &  &  &  &  &  &  &  &  &  \\
RF-02 &  & X &  &  &  &  &  &  &  &  &  &  &  &  &  &  &  &  \\
RF-03 &  &  & X &  &  &  &  &  &  &  & X &  &  &  &  &  &  &  \\
RF-04 &  &  &  & X &  &  &  &  &  &  &  &  &  &  &  &  &  &  \\
RF-05 &  &  &  &  & X &  &  &  &  &  &  &  &  &  &  &  &  &  \\
RF-06 &  &  &  &  &  & X &  &  &  &  &  &  &  &  &  &  &  &  \\
RF-07 &  &  &  &  &  &  & X &  &  &  &  &  &  &  &  &  &  &  \\
RF-08 &  &  &  &  &  &  &  & X &  &  &  &  &  &  &  &  &  &  \\
RF-09 &  &  &  &  &  &  &  &  & X &  &  &  &  &  &  &  &  &  \\
RF-10 &  &  &  &  &  &  &  &  &  & X &  &  &  &  &  &  &  &  \\
RF-11 &  &  &  &  &  &  &  &  &  &  & X &  &  &  &  &  &  &  \\
RF-12 &  &  &  &  &  &  &  &  &  &  &  & X &  &  &  &  &  &  \\
RF-13 &  &  &  &  &  &  &  &  &  &  &  &  & X &  &  &  &  &  \\
RF-14 &  &  &  &  &  &  &  &  &  &  &  &  &  & X &  &  &  &  \\
RF-15 &  &  &  &  &  &  &  &  &  &  &  &  &  &  & X &  &  &  \\
RF-16 &  &  &  &  &  &  &  &  &  &  &  &  &  &  &  & X &  &  \\
RF-17 &  &  &  &  &  &  &  &  &  &  &  &  &  &  &  &  & X &  \\
RF-18 &  &  &  &  &  &  &  &  &  &  &  &  &  &  &  &  &  & X \\
\midrule
RNF-01 &  &  & X & X &  & X & X &  &  & X &  &  &  &  &  &  &  &  \\
RNF-02 & X & X & X & X & X & X & X & X & X & X & X & X & X & X & X & X & X & X \\
RNF-03 &  &  & X &  &  & X &  &  &  &  &  &  &  &  &  &  &  &  \\
RNF-04 &  &  &  &  &  & X & X &  &  & X &  &  &  &  &  &  &  &  \\
RNF-05 &  &  & X &  &  &  &  &  &  &  & X &  &  &  &  &  &  &  \\
RNF-06 &  &  & X &  &  &  &  &  &  &  & X &  &  &  &  &  &  &  \\
RNF-07 &  &  &  &  &  &  &  &  &  &  &  &  &  & X &  &  &  &  \\
RNF-08 &  &  &  &  &  &  &  &  &  &  &  &  & X &  &  &  &  & X \\
\bottomrule
\end{tabular}
\end{table}

\clearpage
\section{Diseño del Sistema}
\label{sec:diseno-sistema}

Esta sección describe el diseño del prototipo AeroTrack, la consola con la que se configura y
opera el sistema. El apartado~\ref{subsec:arquitectura} presenta los seis componentes que
organizan el procesamiento de las imágenes y el motor de identidad entre cámaras. Los
apartados siguientes documentan la consola: la gestión de usuarios por roles, la gestión de
proyectos independientes, el asistente de configuración y las alertas de seguridad con aviso
por correo.

\subsection{Arquitectura propuesta}
\label{subsec:arquitectura}

El prototipo organiza el procesamiento en seis componentes, cada uno con una
responsabilidad única (Figura~\ref{fig:arquitectura}). Todo se ejecuta en el equipo local: el
servidor solo acepta conexiones del propio equipo y ningún fotograma sale de él.

\begin{figure}[H]
\centering
\resizebox{\anchofigura}{!}{%
\begin{tikzpicture}[
  comp/.style={draw=black,thick,rounded corners=2pt,align=center,font=\small,
               minimum width=7.6cm,minimum height=1.25cm,inner sep=4pt},
  dato/.style={draw=black,dashed,rounded corners=2pt,align=center,font=\small,
               minimum width=4.6cm,minimum height=1.1cm,inner sep=4pt},
  flecha/.style={-{Stealth[length=2.2mm]},thick}]
  \node[comp] (a) at (0,0) {\textbf{1. Adquisición}\\archivo de video, cámara USB o flujo de red, por cámara};
  \node[comp,below=0.55cm of a] (d) {\textbf{2. Detección}\\YOLO11 (variante según hardware) y filtros de\\zona útil, tamaño relativo e inmovilidad};
  \node[comp,below=0.55cm of d] (s) {\textbf{3. Seguimiento por cámara}\\ByteTrack en dos etapas sobre el punto de los pies,\\con caja predicha y firma de color del torso};
  \node[comp,below=0.55cm of s] (c) {\textbf{4. Calibración y proyección}\\homografía por cámara hacia el plano (opcional)};
  \node[comp,below=0.55cm of c] (i) {\textbf{5. Identidad entre cámaras}\\vector OSNet de cuerpo entero, tramos, compuerta de\\tiempo y plano, reagrupación al cerrar la sesión};
  \node[comp,below=0.55cm of i] (p) {\textbf{6. Persistencia, analítica y presentación}\\conteos, alertas, reportes y consola web};
  \node[dato,right=1.0cm of i] (m) {Memoria de apariencia\\(base local, 24 h por defecto)};
  \node[dato,right=1.0cm of p] (g) {Grabación de la sesión\\(identificadores y posiciones)};
  \foreach \x/\y in {a/d,d/s,s/c,c/i,i/p} { \draw[flecha] (\x) -- (\y); }
  \draw[flecha,{Stealth[length=2.2mm]}-{Stealth[length=2.2mm]}] (i) -- (m);
  \draw[flecha] (p) -- (g);
\end{tikzpicture}}
\caption{Arquitectura del prototipo AeroTrack. Los recuadros discontinuos son los datos por
persona que el sistema conserva en el equipo.}
\label{fig:arquitectura}
\end{figure}

El primer componente es la adquisición, que lee cada fuente de cámara, sea un archivo de
video, una cámara USB o un flujo de red, y la abre una sola vez. El segundo es la detección: el
detector YOLO11 (Jocher y Qiu, 2024) localiza personas en cada muestra con una caja y una
confianza. La variante del detector no la elige el operador, sino el sistema según el hardware
del equipo: nano en CPU, y pequeña o mediana con GPU. Tras la detección se aplican tres filtros
contra falsos positivos: se descarta lo que cae fuera de la zona útil de la imagen o del área de
trabajo del plano, las detecciones dudosas mucho más pequeñas que las personas seguras de la
misma cámara (gorros, carteles) y los objetos que nunca se mueven y se detectan con baja
confianza.

El tercero es el seguimiento dentro de cada cámara. Aplica la asociación en dos etapas de
ByteTrack (Zhang et al., 2022) sobre el punto de los pies, que es el que la homografía puede
proyectar, con un filtro de Kalman de velocidad constante: primero se emparejan las
detecciones de alta confianza y después las de baja confianza, para no perder a una persona
parcialmente ocluida. Dentro de la puerta espacial, las coincidencias se ordenan además con
dos señales al estilo de BoT-SORT (Aharon et al., 2022): la superposición entre la caja predicha
y la detectada, y una firma de color de la ropa calculada sobre dos franjas del torso.

El cuarto componente es la calibración y proyección: cada cámara tiene una homografía propia
hacia el plano, calculada con al menos cuatro correspondencias entre el suelo del video y el
plano, que se valida por consistencia geométrica antes de aceptarse. La calibración es
opcional; sin ella la cámara cuenta y sigue personas, pero no las ubica en el plano.

El quinto es el motor de identidad entre cámaras (Figura~\ref{fig:identidad}). Para cada vista
confiable de una persona calcula un vector de apariencia OSNet (Zhou et al., 2019) de 512
dimensiones sobre el recorte de cuerpo entero, con las partes tapadas por otra persona
rellenadas de gris, y lo acumula en el tramo continuo de seguimiento al que pertenece. Una
persona lleva un identificador provisional hasta reunir suficientes vistas y solo entonces recibe
un identificador confirmado, que es el único que cuenta en los conteos. La unión entre cámaras
exige a la vez parecido de apariencia y coherencia de tiempo y, si hay calibración, de posición
en el plano, y solo se intenta entre cámaras que el operador declaró relacionadas y cuyos relojes
declaró sincronizados; ante una candidata ambigua, la asociación se marca como incierta en
lugar de forzarse. Con la memoria de apariencia activa, que es la configuración por defecto, el
motor compara además cada identidad nueva con todas las personas guardadas en una base
de datos local del proyecto durante el plazo de retención (24 horas por defecto, hasta 168), de
modo que reconoce a la misma persona en sesiones y días posteriores. Al cerrar la sesión,
reagrupa los tramos con la grabación completa a la vista y renumera a las personas de 1 a N.

\begin{figure}[H]
\centering
\resizebox{\anchofigura}{!}{%
\begin{tikzpicture}[
  paso/.style={draw=black,thick,rounded corners=2pt,align=center,font=\small,
               minimum width=6.4cm,minimum height=1.0cm,inner sep=3pt},
  decision/.style={draw=black,thick,diamond,aspect=2.6,align=center,font=\small,inner sep=1pt},
  fin/.style={draw=black,thick,rounded corners=6pt,align=center,font=\small,
              minimum width=3.6cm,minimum height=0.9cm,inner sep=3pt},
  flecha/.style={-{Stealth[length=2.2mm]},thick}]
  \node[paso] (n1) at (0,0) {Vista confiable de una persona seguida\\(confianza, altura, borde, oclusión)};
  \node[paso,below=0.5cm of n1] (n2) {Vector OSNet de cuerpo entero (512 dimensiones)\\y prototipo del tramo de seguimiento};
  \node[decision,below=0.5cm of n2] (d1) {¿Vistas suficientes\\para confirmar?};
  \node[fin,left=1.2cm of d1] (f0) {Identificador\\provisional};
  \node[decision,below=0.6cm of d1] (d2) {¿Coincide con una\\persona guardada?};
  \node[fin,right=1.2cm of d2] (f1) {Recupera su ID\\(reidentificada)};
  \node[decision,below=0.6cm of d2] (d3) {¿Candidata en cámara\\relacionada, coherente en\\apariencia, tiempo y plano?};
  \node[fin,left=1.2cm of d3] (f3) {ID nuevo\\(local)};
  \node[decision,below=0.6cm of d3] (d4) {¿Es inequívoca?};
  \node[fin,right=1.2cm of d4] (f2) {Reutiliza su ID\\(estimada)};
  \node[fin,left=1.2cm of d4] (f4) {No fuerza la unión\\(incierta)};
  \draw[flecha] (n1) -- (n2);
  \draw[flecha] (n2) -- (d1);
  \draw[flecha] (d1) -- node[above,font=\small]{no} (f0);
  \draw[flecha] (d1) -- node[right,font=\small]{sí} (d2);
  \draw[flecha] (d2) -- node[above,font=\small]{sí} (f1);
  \draw[flecha] (d2) -- node[right,font=\small]{no} (d3);
  \draw[flecha] (d3) -- node[above,font=\small]{no} (f3);
  \draw[flecha] (d3) -- node[right,font=\small]{sí} (d4);
  \draw[flecha] (d4) -- node[above,font=\small]{sí} (f2);
  \draw[flecha] (d4) -- node[above,font=\small]{no} (f4);
  \node[font=\small\itshape,align=center,below=0.5cm of d4]
    {La consulta a personas guardadas solo ocurre con la memoria de apariencia activa,\\que es la configuración por defecto (RNF-01 no conforme).};
\end{tikzpicture}}
\caption{Decisión de identidad de una persona en el motor de identidad entre cámaras.}
\label{fig:identidad}
\end{figure}

El sexto componente reúne persistencia, analítica y presentación. Calcula la ocupación por
zona, el flujo por los accesos y las alertas de aglomeración a partir de las personas
confirmadas, y lo expone en la consola, que ofrece configuración, plano en vivo, video por
cámara, alertas y reportes. No se limita a mantener un estado agregado en memoria: guarda en
el equipo una grabación de cada sesión con, por muestra y por persona, el identificador, la caja,
la posición en el plano y su trayectoria reciente, sin plazo de retención automático, y a partir de
esas grabaciones calcula las métricas comerciales por negocio. No guarda fotogramas de las
cámaras en vivo ni vectores de apariencia en la grabación. Las consecuencias de esa
persistencia se examinan en el Anexo~\ref{anx:marco}, apartado B.9.

\subsection{Gestión de usuarios y roles}
\label{subsec:diseno-usuarios}

El acceso a la consola requiere usuario y contraseña. En la primera ejecución el sistema pide
crear un primer usuario (Figura~\ref{fig:primer-usuario}), que queda con el rol de operador y
puede dar de alta a los demás. La contraseña debe tener al menos seis caracteres y se guarda
en el equipo como resumen criptográfico con sal propia, nunca en texto plano.

\figura{aerotrack-01-primer-usuario.png}
  {Creación del primer usuario en la primera ejecución de la consola.}
  {fig:primer-usuario}

Una vez dentro, la vista general muestra en la cabecera el proyecto activo junto al usuario y su
rol, y organiza la consola en pestañas: vista general, alertas, videos y resultados, reportes y
configuración (Figura~\ref{fig:vista-general}). Desde esta vista se filtra por nivel, cámara, zona
de imagen y rendimiento, y se inicia el monitoreo sobre el plano del nivel. En la captura los
contadores de cámaras seleccionadas, entradas, salidas y tiempo de análisis están en cero
porque el proyecto aún no tiene cámaras configuradas.

\figura{aerotrack-02-vista-general.png}
  {Vista general de la consola, con el proyecto activo en la cabecera y el plano del nivel 3.}
  {fig:vista-general}

El sistema distingue dos roles (Figura~\ref{fig:usuarios}). El operador configura cámaras y
plano; el administrador consulta resultados y ajusta los umbrales de alerta, sin acceso a la
configuración. Los roles se aplican en el servidor y no solo en la interfaz. El panel de usuarios
lista las cuentas existentes con su rol, permite eliminarlas y permite crear otras indicando
nombre, contraseña y rol.

\figura{aerotrack-03-usuarios.png}
  {Panel de usuarios del sistema y formulario de alta con selección de rol.}
  {fig:usuarios}

\subsection{Gestión de proyectos}
\label{subsec:diseno-proyectos}

Un proyecto representa un espacio monitoreado, por ejemplo un nivel de un terminal. Cada
proyecto guarda por separado su propio plano, sus cámaras, su calibración, sus incidentes y su
memoria de apariencia, de modo que la configuración y los datos de un espacio no se mezclan
con los de otro. Al entrar en la configuración, la consola ofrece dos acciones: crear un proyecto
nuevo, que conduce al asistente descrito en el apartado~\ref{subsec:diseno-asistente}, o añadir
una cámara al proyecto abierto (Figura~\ref{fig:crear-proyecto}). Para crear el proyecto basta
con darle un nombre.

\figura{aerotrack-04-crear-proyecto.png}
  {Opciones para crear un proyecto nuevo o añadir una cámara al proyecto abierto.}
  {fig:crear-proyecto}

Los proyectos guardados se listan con su nombre, el espacio al que corresponden, el número
de cámaras y de planos, el momento de la última edición y su estado
(Figura~\ref{fig:proyectos-guardados}). En la captura figuran dos proyectos, ambos marcados
como incompletos: el proyecto principal, sobre el nivel 3 del aeropuerto y con cuatro planos, y
el proyecto Terminal~A, Nivel~1, con un plano. Desde la lista se abre un proyecto, se retoma su
configuración en el punto en que quedó, se renombra directamente sobre su tarjeta
(Figura~\ref{fig:renombrar}) o se elimina.

\figura{aerotrack-05-proyectos-guardados.png}
  {Lista de proyectos guardados, cada uno con sus cámaras, planos y estado.}
  {fig:proyectos-guardados}

\figura{aerotrack-06-renombrar-proyecto.png}
  {Cambio de nombre de un proyecto sobre su propia tarjeta.}
  {fig:renombrar}

Al entrar en la consola, una pantalla pregunta en qué proyecto se quiere trabajar, muestra para
cada uno su espacio, su número de cámaras y el estado de su configuración, y señala el último
abierto (Figura~\ref{fig:selector-proyecto}). Así el usuario elige de forma expresa el espacio en
el que va a trabajar y no mezcla datos de dos lugares distintos.

\figura{aerotrack-07-selector-proyecto.png}
  {Selección del proyecto de trabajo al entrar en la consola.}
  {fig:selector-proyecto}

\subsection{Asistente de configuración}
\label{subsec:diseno-asistente}

La configuración de un proyecto se completa con un asistente. Las capturas de este apartado
corresponden a su versión de ocho pasos: proyecto, plano, conexión, prueba, ubicación,
calibración, zonas y alertas, y revisión. Una barra superior muestra los pasos, marca los
completados y lleva la cuenta de cuántos están listos, y una indicación en la cabecera confirma
que el avance queda guardado, de modo que un proyecto incompleto puede retomarse desde la
lista de proyectos. Las capturas de los pasos 4, 6 y 8 se tomaron sin video, porque no se
dispone de grabaciones del terminal.

La versión vigente del asistente conserva las mismas tareas, pero las reúne en cinco etapas. La
Tabla~\ref{tab:asistente} muestra la correspondencia, y los pasos se describen a continuación
en el orden de las capturas.

\begin{table}[H]
\centering
\small
\caption{Correspondencia entre los pasos de las capturas y las etapas del asistente vigente.}
\label{tab:asistente}
\begin{tabular}{@{}L{5.2cm}L{9.4cm}@{}}
\toprule
\textbf{Etapa vigente} & \textbf{Pasos de las capturas que reúne} \\
\midrule
Proyecto & Paso 1 (datos del espacio y tipo de fuente) y paso 2 (plano) \\
Cámaras & Paso 3 (cámaras, fuente, zona útil y accesos) y paso 4 (prueba y vista en vivo) \\
Homografía y cámaras relacionadas & Paso 5 (ubicación) y paso 6 (calibración, ahora opcional), más la relación entre cámaras y la sincronización de sus relojes \\
Contexto comercial & Paso 7 (zonas, negocios y parámetros de aviso de aglomeración) \\
Validación & Paso 8 (revisión) \\
\bottomrule
\end{tabular}
\end{table}
\verificar{ubicación de la cámara sobre el plano (paso 5 de las capturas) dentro de la etapa de
homografía del asistente vigente: inferida de la estructura de la consola, sin captura nueva que
la muestre.}

\subsubsection*{Paso 1: proyecto}

Se identifica el espacio con el nombre del aeropuerto o proyecto y el del terminal y nivel, y se
elige el tipo de fuente: videos grabados, para analizar grabaciones y consultar resultados por
periodo, o cámaras en vivo, conectadas por USB, RTSP o IP (Figura~\ref{fig:paso1}). Una nota
lateral recuerda que el análisis es de presencia y recorridos, sin reconocimiento facial.

\figura{aerotrack-08-paso1-proyecto.png}
  {Paso 1 del asistente: datos del espacio y tipo de fuente.}
  {fig:paso1}

\subsubsection*{Paso 2: plano}

Se define el plano del espacio: se elige uno de los planos del aeropuerto, se importa una
imagen o un PDF, o se dibuja el espacio con rectángulos, polígonos y líneas. El plano se
acompaña de su ancho y su alto en metros, y admite marcar un área de trabajo como
referencia opcional (Figura~\ref{fig:paso2}).

\figura{aerotrack-09-paso2-plano.png}
  {Paso 2 del asistente: plano del nivel y sus dimensiones en metros.}
  {fig:paso2}

\subsubsection*{Paso 3: conexión}

Se añaden las cámaras del proyecto (Figura~\ref{fig:paso3}). Para cada una se indica el plano
asociado, el nombre visible en el mapa, la ubicación, la orientación y la fuente, con la ruta del
video o la dirección de la cámara, y se fijan el umbral de alerta en personas y el tiempo de
concentración sostenida (Figura~\ref{fig:paso3-camara}). En el mismo paso se marca la zona
útil de la imagen, que deja fuera reflejos, otras plantas y áreas que no deben analizarse, y, si la
cámara ve una puerta, un acceso formado por un área exterior y otra interior, entre las que se
cuenta el paso de personas (Figura~\ref{fig:paso3-zona}).

\figura{aerotrack-10-paso3-conexion.png}
  {Paso 3 del asistente: lista de cámaras del proyecto, todavía vacía.}
  {fig:paso3}

\figura{aerotrack-11-paso3-camara.png}
  {Paso 3 del asistente: datos de la cámara, fuente y umbrales de alerta.}
  {fig:paso3-camara}

\figura{aerotrack-12-paso3-zona-util.png}
  {Paso 3 del asistente: zona útil de la imagen y acceso de un local.}
  {fig:paso3-zona}

\subsubsection*{Paso 4: prueba}

Se comprueba la cámara en movimiento para revisar encuadre, luz y retardo. El panel muestra
el estado de la fuente, su resolución, su frecuencia de cuadros y si está calibrada, y permite
iniciar una prueba de seguimiento en modo laboratorio que presenta el video procesado con
identificadores temporales (Figura~\ref{fig:paso4}).

\figura{aerotrack-13-paso4-prueba.png}
  {Paso 4 del asistente: prueba de la cámara, con el visor sin señal.}
  {fig:paso4}

\subsubsection*{Paso 5: ubicación}

Se coloca la cámara sobre el plano en el punto donde está instalada y se orienta hacia donde
apunta, con su dirección, alcance, altura y ancho de cobertura (Figura~\ref{fig:paso5}). Esa
cobertura es orientativa: la precisión la aporta la calibración del paso siguiente.

\figura{aerotrack-14-paso5-ubicacion.png}
  {Paso 5 del asistente: posición y cobertura de la cámara sobre el plano.}
  {fig:paso5}

\subsubsection*{Paso 6: calibración}

Se relaciona el suelo del video con el plano marcando puntos de referencia, como esquinas o
cruces de baldosas, primero en la imagen y después en el mapa (Figura~\ref{fig:paso6}). Se
requieren al menos cuatro puntos repartidos; con cinco o más, cada referencia se comprueba
además dejándola fuera del ajuste. En la versión vigente la calibración es opcional, y la misma
etapa recoge qué cámaras están relacionadas y si sus relojes están sincronizados, condición sin
la cual el sistema no asocia personas entre cámaras.

\figura{aerotrack-15-paso6-calibracion.png}
  {Paso 6 del asistente: calibración del suelo del video con el plano.}
  {fig:paso6}

\subsubsection*{Paso 7: zonas y alertas}

Se dibujan sobre el plano las zonas de interés con polígonos, rectángulos o círculos, y se define
cuándo avisar de una aglomeración mediante un radio en metros, un número mínimo de
personas y un tiempo de permanencia (Figura~\ref{fig:paso7}). La pantalla advierte que una fila y
una aglomeración pueden parecerse geométricamente, por lo que los límites deben ajustarse al
uso de cada zona. En la versión vigente esta etapa recoge además los negocios del piso, que
alimentan el panel comercial.

\figura{aerotrack-16-paso7-zonas.png}
  {Paso 7 del asistente: trazado de una zona y parámetros de aviso de aglomeración.}
  {fig:paso7}

\subsubsection*{Paso 8: revisión}

Se revisa el proyecto antes de abrir el monitoreo: una lista indica qué está preparado y qué
sigue pendiente, y un resumen muestra el número de cámaras, zonas dibujadas y fuentes
probadas, junto con la escala (Figura~\ref{fig:paso8}). En la captura, con la cámara sin probar ni
calibrar, la opción de guardar y abrir el monitoreo aparece desactivada y queda disponible la
de continuar las pruebas en laboratorio. La pantalla aclara que una verificación completada no
sustituye la validación de precisión.

\figura{aerotrack-17-paso8-revision.png}
  {Paso 8 del asistente: revisión del proyecto antes de abrir el monitoreo.}
  {fig:paso8}

\subsection{Alertas de seguridad con aviso por correo}
\label{subsec:diseno-alertas}

La vista de alertas reúne lo que necesita el personal de seguridad para atender una
aglomeración (Figura~\ref{fig:alertas-correo}). El apartado \enquote{Cuándo avisar} fija los
umbrales que definen una aglomeración: el número mínimo de personas, la permanencia en
segundos y el radio en metros; en la captura, 4 personas, 3 segundos y 1,5 metros. Estos
umbrales se cambian desde la misma vista, sin entrar en la configuración, y son el único ajuste
que también puede hacer el administrador.

El panel \enquote{Avisar por correo}, desplegado en la captura y todavía desactivado, pide la
cuenta que envía, su contraseña de aplicación y los correos que reciben el aviso, separados por
coma, hasta un máximo de veinte, y permite enviar un correo de prueba. Una nota aclara que el
servicio de correo no acepta la contraseña normal de la cuenta, sino una contraseña de
aplicación generada con la verificación en dos pasos. Debajo, tres contadores muestran las
alertas activas, el total de la sesión y los incidentes sin revisar, y la bitácora de incidentes
conserva cada caso aunque la sesión termine, para que el personal lo marque como revisado,
falsa alarma o resuelto. En la captura la bitácora aparece sin registros porque no se han
procesado grabaciones del terminal.

Con el aviso activado, el sistema envía un correo por cada incidente nuevo de aglomeración,
con el lugar, la cámara y el motivo del aviso, y la advertencia de que requiere comprobación
humana; un aviso ya enviado no se reenvía y una misma zona no genera más de un correo cada
120 segundos. La contraseña se guarda solo en el equipo, fuera de los archivos del proyecto: la
interfaz nunca la muestra de vuelta y no se copia al duplicar un proyecto.

\figura{aerotrack-18-alertas-correo.png}
  {Vista de alertas: umbrales de aglomeración, aviso por correo desplegado, contadores y
   bitácora de incidentes.}
  {fig:alertas-correo}

\clearpage
\section{Planificación y Gestión del Proyecto}
\label{sec:planificacion}

\textbf{Cronograma y Tareas.} El proyecto se planifica sobre el calendario del curso, con
cuatro presentaciones quincenales de avance (semanas 5, 8, 11 y 14) y la sustentación final en
la semana 16.

\begin{center}
\begin{tabular}{@{}ll@{}}
\toprule
\textbf{Estado} & \\
\midrule
Completado & \\
En curso & \\
Pendiente & \\
Programado & \\
\bottomrule
\end{tabular}
\end{center}

\figuraPendiente{gantt.png}{\textwidth}{10cm}{Cronograma de fases y actividades del
proyecto Aerotrack, semanas S1 a S16.}{fig:gantt}

Las fases de organización del equipo, antecedentes y desarrollo del prototipo (incluida la
incorporación posterior del monitoreo integrado) están completadas. Requisitos y arquitectura
tiene completos los requisitos funcionales y los doce casos de uso, pero mantiene dos
decisiones abiertas: la arquitectura definitiva (granularidad de salida) y la localización del
procesamiento. Licencias de datos públicos y pruebas/validación están en curso: las pruebas
sintéticas y la validación cualitativa sobre videos propios están completas; la solicitud formal
de acceso a datos del terminal sigue pendiente de respuesta de LAP. La redacción del informe
está en curso, de la cual este documento es la entrega más reciente. La entrega y sustentación
final es la única fase íntegramente pendiente, programada para las semanas 15 y 16.

\subsection{Roles y Responsabilidades}
\label{subsec:roles-responsabilidades}

\begin{longtable}{@{}L{4cm}L{4cm}L{7.5cm}@{}}
\toprule
\textbf{Integrante} & \textbf{Frente} & \textbf{Responsabilidad principal} \\
\midrule
\endhead
Angel Aguilar & Liderazgo e integración & Coordinación general; consolidación de secciones;
  redacción de este documento \\
Stephano Rivadeneyra & Ciencia de datos & Metodología del informe oficial; desarrollo del
  monitoreo integrado (conteo, seguimiento por flujo, cartografía, dashboard) \\
José Ortega & Ingeniería de datos & Resultados del informe oficial; verificación de licencias
  de los conjuntos de datos públicos \\
Fabián Moreno Ugarte & Validación ética, legal y documentación & Antecedentes, marco
  normativo, análisis de impacto y anexo del informe oficial \\
\bottomrule
\caption{Roles y responsabilidades del equipo Aerotrack.}
\label{tab:roles}
\end{longtable}

\clearpage
\section{Referencias Bibliográficas}
\label{sec:referencias}

% Formato APA con sangria francesa (hanging indent), tal como en el
% informe original. No se usa biblatex/biber: las citas dentro del
% texto van escritas literalmente como texto parentetico, igual que en
% el documento fuente (p. ej. "(Li, Zhang y Chen, 2018)").
\newcommand{\refapa}[1]{\par\noindent\hangindent=1.25cm\hangafter=1 #1\par\vspace{0.35cm}}

% \url no puede escribirse directamente dentro del argumento de \refapa cuando
% contiene caracteres especiales de TeX (%, _, #, ~, ^, &): ese argumento se
% tokeniza ANTES de que \url cambie los codigos de categoria, asi que "_" abre
% subindice y "%" abre un comentario que se traga el resto de la linea (el
% error "File ended while scanning use of \refapa" viene de ahi). \urldef
% arma el enlace por separado, en la categoria de codigos correcta, y aqui
% solo se referencia el nombre ya seguro.
\urldef{\urlConstitucion}\url{https://leyes.congreso.gob.pe/Documentos/constituciones_ordenado/CONSTIT_1993/Texto_actualizado_CONS_1993.pdf}

\refapa{Congreso de la República del Perú. (1993). \emph{Constitución Política del Perú},
artículo 2. Texto actualizado: \urlConstitucion}

\refapa{Congreso de la República del Perú. (2011). \emph{Ley N.º 29733, Ley de Protección
de Datos Personales}. Diario Oficial El Peruano.
\url{https://www.gob.pe/institucion/congreso-de-la-republica/normas-legales/243470-29733}}

\refapa{Congreso de la República del Perú. (2013). \emph{Ley N.º 30120, Ley de apoyo a la
seguridad ciudadana con cámaras de videovigilancia públicas y privadas}. Diario Oficial El
Peruano. \url{https://www.leyes.congreso.gob.pe/Documentos/Leyes/30120.pdf}}

\urldef{\urlDirectiva}\url{https://www.minjus.gob.pe/wp-content/uploads/2020/01/Directiva-N°-01-2020-DGTAIPD-1.pdf}

\refapa{Dirección General de Transparencia, Acceso a la Información Pública y Protección
de Datos Personales. (2020). \emph{Directiva N.º 01-2020-JUS/DGTAIPD, para el tratamiento
de datos personales mediante sistemas de videovigilancia}. Ministerio de Justicia y Derechos
Humanos. \urlDirectiva}

\refapa{Ministerio de Justicia y Derechos Humanos. (2024). \emph{Decreto Supremo N.º
016-2024-JUS, que aprueba el Reglamento de la Ley N.º 29733, Ley de Protección de Datos
Personales}. Diario Oficial El Peruano.
\url{https://www.gob.pe/institucion/smv/normas-legales/6426760-016-2024-jus}}

\refapa{Ministerio del Interior. (2020). \emph{Decreto Supremo N.º 007-2020-IN, que
aprueba el Reglamento del Decreto Legislativo N.º 1218 y de la Ley N.º 30120}. Diario Oficial
El Peruano.
\url{https://www.gob.pe/institucion/mininter/normas-legales/1081739-007-2020-in}}

\refapa{Poder Ejecutivo. (2015). \emph{Decreto Legislativo N.º 1218, que regula el uso de las
cámaras de videovigilancia}. Diario Oficial El Peruano.
\url{https://busquedas.elperuano.pe/normaslegales/decreto-legislativo-que-regula-el-uso-de-las-camaras-de-vide-decreto-legislativo-n-1218-1291565-8/}}

\bigskip
\noindent\textbf{Trabajos técnicos}
\medskip

\refapa{Aharon, N., Orfaig, R., \& Bobrovsky, B.-Z. (2022). BoT-SORT: Robust associations
multi-pedestrian tracking. \emph{arXiv preprint arXiv:2206.14651}.
\url{https://arxiv.org/abs/2206.14651}}

\refapa{Dalal, N., \& Triggs, B. (2005). Histograms of oriented gradients for human detection.
En \emph{Proceedings of the IEEE Computer Society Conference on Computer Vision and Pattern
Recognition} (pp. 886-893). IEEE. \url{https://doi.org/10.1109/CVPR.2005.177}}

\refapa{Dendorfer, P., Rezatofighi, H., Milan, A., Shi, J., Cremers, D., Reid, I., Roth, S.,
Schindler, K., \& Leal-Taixé, L. (2020). MOT20: A benchmark for multi object tracking in crowded
scenes. \emph{arXiv preprint arXiv:2003.09003}. \url{https://arxiv.org/abs/2003.09003}}

\refapa{Jocher, G., \& Qiu, J. (2024). \emph{Ultralytics YOLO11} (versión 11.0.0) [Software].
Ultralytics. \url{https://github.com/ultralytics/ultralytics}}

\refapa{Li, Y., Zhang, X., \& Chen, D. (2018). CSRNet: Dilated convolutional neural
networks for understanding the highly congested scenes. En \emph{Proceedings of the IEEE
Conference on Computer Vision and Pattern Recognition} (pp. 1091-1100). IEEE.
\url{https://arxiv.org/abs/1802.10062}}

\refapa{Lin, W., Zhao, C., \& Chan, A. B. (2025). Point-to-region loss for semi-supervised
point-based crowd counting. En \emph{Proceedings of the IEEE/CVF Conference on Computer
Vision and Pattern Recognition} (pp. 29363-29373). IEEE.
\url{https://arxiv.org/abs/2505.21943}}

\refapa{Redmon, J., Divvala, S., Girshick, R., \& Farhadi, A. (2016). You only look once:
Unified, real-time object detection. En \emph{Proceedings of the IEEE Conference on Computer
Vision and Pattern Recognition} (pp. 779-788). IEEE. \url{https://arxiv.org/abs/1506.02640}}

\refapa{Song, Q., Wang, C., Jiang, Z., Wang, Y., Tai, Y., Wang, C., Li, J., Huang, F., \& Wu,
Y. (2021). Rethinking counting and localization in crowds: A purely point-based framework.
En \emph{Proceedings of the IEEE/CVF International Conference on Computer Vision}
(pp. 3365-3374). IEEE. \url{https://arxiv.org/abs/2107.12746}}

\refapa{Zhang, Y., Sun, P., Jiang, Y., Yu, D., Weng, F., Yuan, Z., Luo, P., Liu, W., \& Wang,
X. (2022). ByteTrack: Multi-object tracking by associating every detection box. En
\emph{Proceedings of the European Conference on Computer Vision} (pp. 1-21). Springer.
\url{https://arxiv.org/abs/2110.06864}}

\refapa{Zhou, K., Yang, Y., Cavallaro, A., \& Xiang, T. (2019). Omni-scale feature learning
for person re-identification. En \emph{Proceedings of the IEEE/CVF International Conference
on Computer Vision} (pp. 3702-3712). IEEE. \url{https://arxiv.org/abs/1905.00953}}

\clearpage
\section{Anexos}
\label{sec:anexos}

% El PDF fuente termina en el encabezado "12. Anexos" sin contenido
% posterior visible. Se deja la seccion abierta para que el contenido
% de los anexos (si existe en una version mas completa del informe) se
% incorpore aqui.


\end{document}
